\chapter{§1.6 小数}
\label{sec:1.6}


\prerequisite{§1.2 四则运算，§1.5 分数（初步）}

\gradelevel{3-5 年级}


\section*{小数的意义}


\begin{explore}


文具店里一支铅笔标价 $1.50$ 元，一块橡皮标价 $0.80$ 元。这些价格中的小数点是什么意思？ $1.50$ 和 $1$ 元 $5$ 角有什么关系？

\end{explore}

\begin{understand}


\textbf{小数}是十进分数的另一种写法。

\[
\frac{1}{10} = 0.1, \quad \frac{1}{100} = 0.01, \quad \frac{1}{1000} = 0.001
\]


小数点左边是整数部分，右边是小数部分。小数点右边第一位是\textbf{十分位}，第二位是\textbf{百分位}，第三位是\textbf{千分位}。

例如 $3.75$ ：
\begin{itemize}
  \item 整数部分： $3$ 。
  \item 十分位： $7$ ，表示 $7$ 个 $0.1$ 。
  \item 百分位： $5$ ，表示 $5$ 个 $0.01$ 。
\end{itemize}


\[
3.75 = 3 + \frac{7}{10} + \frac{5}{100} = 3 + 0.7 + 0.05 = \frac{375}{100}
\]


小数与分数的关系：小数就是分母为 $10$ 、 $100$ 、 $1000$ ……的分数。

\end{understand}

\begin{workedexamples}


\textbf{例 1.} 把 $\frac{3}{10}$ 、 $\frac{47}{100}$ 写成小数。

\textbf{解：} $\frac{3}{10} = 0.3$ ， $\frac{47}{100} = 0.47$ 。

\textbf{例 2.} 把 $0.125$ 写成分数。

\textbf{解：} $0.125 = \frac{125}{1000} = \frac{1}{8}$ （约分： $125$ 和 $1000$ 的最大公因数是 $125$ ）。

\end{workedexamples}

\begin{keytakeaway}


\begin{itemize}
  \item 小数是分母为 $10$ 、 $100$ 、 $1000$ 等的分数的简写。
  \item 小数点右边每一位的意义：十分位、百分位、千分位……
  \item 小数和分数可以互相转化。
\end{itemize}


\end{keytakeaway}

\begin{practice}


\begin{enumerate}
  \item 写出 $2.36$ 中各数位上的数字及其意义。
  \item 把 $\frac{9}{100}$ 和 $\frac{23}{1000}$ 写成小数。
  \item 把 $0.75$ 写成最简分数。
\end{enumerate}


\end{practice}

\section*{小数的读写与大小比较}


\begin{explore}


超市里苹果每千克 $5.60$ 元，梨每千克 $5.08$ 元。哪种水果更贵？ $5.60$ 和 $5.08$ 应该怎么比较？

\end{explore}

\begin{understand}


\textbf{小数的大小比较}方法：
\begin{enumerate}
  \item 先比较整数部分，整数部分大的小数大。
  \item 整数部分相同时，从十分位起逐位比较。
\end{enumerate}


例如比较 $5.60$ 和 $5.08$ ：整数部分都是 $5$ ，十分位 $6 > 0$ ，所以 $5.60 > 5.08$ 。

\textbf{小数的性质}：小数末尾添上或去掉零，小数的大小不变。 $5.60 = 5.6$ 。

\end{understand}

\begin{workedexamples}


\textbf{例 1.} 把 $3.5$ 、 $3.05$ 、 $3.50$ 、 $3.55$ 从小到大排列。

\textbf{解：}
\begin{itemize}
  \item $3.5 = 3.50$ （末尾添零）。
  \item 比较： $3.05 < 3.50 = 3.5 < 3.55$ 。
\end{itemize}


从小到大： $3.05 < 3.5 = 3.50 < 3.55$ 。

\textbf{例 2.} 在 $\bigcirc$ 中填 $>$ 、 $<$ 或 $=$ ： $0.9 \enspace\bigcirc\enspace 0.85$ 。

\textbf{解：} 整数部分都是 $0$ 。十分位： $9 > 8$ ，所以 $0.9 > 0.85$ 。

\end{workedexamples}

\begin{keytakeaway}


\begin{itemize}
  \item 比较小数从高位到低位逐位比较。
  \item 小数末尾的零不影响大小。
  \item 比较时可以补零对齐位数再比较。
\end{itemize}


\end{keytakeaway}

\begin{practice}


\begin{enumerate}
  \item 比较大小： $4.7 \enspace\underline{\quad}\enspace 4.07$ 。
  \item 把 $0.6$ 、 $0.06$ 、 $0.60$ 、 $0.66$ 从小到大排列。
  \item 在 $2.3$ 和 $2.4$ 之间写出一个小数。
\end{enumerate}


\end{practice}

\section*{小数的四则运算}


\begin{explore}


小明买了一本书 $12.50$ 元和一个笔记本 $3.80$ 元。他一共花了多少钱？需要找回多少钱（如果付了 $20$ 元）？小数该怎么加减？

\end{explore}

\begin{understand}


\textbf{小数加减法}：小数点对齐，按照整数加减法的方法计算，得数的小数点与上面对齐。

\[
\begin{aligned}
&12.50 \\
+\enspace&\phantom{0}3.80 \\
\hline
&16.30
\end{aligned}
\]


\textbf{小数乘法}：先按整数乘法计算（忽略小数点），再数两个因数的小数位数之和，从积的末尾起数出相应位数，点上小数点。

例如 $0.3 \times 0.4$ ： $3 \times 4 = 12$ ，两个因数共有 $2$ 位小数，所以积是 $0.12$ 。

\textbf{小数除法}：
\begin{itemize}
  \item 除数是整数：按整数除法算，商的小数点与被除数的小数点对齐。
  \item 除数是小数：先移动小数点把除数变成整数，被除数的小数点也相应移动。
\end{itemize}


例如 $7.2 \div 0.6$ ：除数和被除数同时乘 $10$ ，变为 $72 \div 6 = 12$ 。

\end{understand}

\begin{workedexamples}


\textbf{例 1.} 计算 $5.73 + 2.8$ 。

\textbf{解：} 小数点对齐， $2.8$ 补零为 $2.80$ ：

\[
\begin{aligned}
&5.73 \\
+\enspace&2.80 \\
\hline
&8.53
\end{aligned}
\]


\textbf{例 2.} 计算 $2.5 \times 0.4$ 。

\textbf{解：} 先算 $25 \times 4 = 100$ 。两个因数共 $2$ 位小数，所以 $2.5 \times 0.4 = 1.00 = 1$ 。

\textbf{例 3.} 计算 $8.4 \div 0.7$ 。

\textbf{解：} 除数和被除数同时乘 $10$ ： $84 \div 7 = 12$ 。所以 $8.4 \div 0.7 = 12$ 。

\end{workedexamples}

\begin{keytakeaway}


\begin{itemize}
  \item 小数加减法：小数点对齐，位数不够补零。
  \item 小数乘法：先按整数乘，再点小数点（小数位数之和）。
  \item 小数除法：把除数变成整数再计算。
\end{itemize}


\end{keytakeaway}

\begin{practice}


\begin{enumerate}
  \item 计算： $6.35 + 4.7$ 。
  \item 计算： $10.2 - 3.85$ 。
  \item 计算： $3.6 \times 0.5$ 。
  \item 计算： $9.6 \div 1.2$ 。
\end{enumerate}


\end{practice}



\begin{answer}


\textbf{练一练 小数的意义}

\begin{enumerate}
  \item 整数部分： $2$ 。十分位： $3$ ，表示 $3$ 个 $0.1$ 。百分位： $6$ ，表示 $6$ 个 $0.01$ 。
  \item $\frac{9}{100} = 0.09$ ， $\frac{23}{1000} = 0.023$ 。
  \item $0.75 = \frac{75}{100} = \frac{3}{4}$ 。
\end{enumerate}


\textbf{练一练 小数的读写与大小比较}

\begin{enumerate}
  \item $4.7 > 4.07$ （十分位 $7 > 0$ ）。
  \item $0.06 < 0.6 = 0.60 < 0.66$ 。
  \item 答案不唯一，例如 $2.35$ （ $2.3$ 和 $2.4$ 之间有无数个小数）。
\end{enumerate}


\textbf{练一练 小数的四则运算}

\begin{enumerate}
  \item $6.35 + 4.7 = 6.35 + 4.70 = 11.05$ 。
  \item $10.2 - 3.85 = 10.20 - 3.85 = 6.35$ 。
  \item $3.6 \times 0.5 = 1.8$ （ $36 \times 5 = 180$ ， $2$ 位小数）。
  \item $9.6 \div 1.2 = 96 \div 12 = 8$ 。
\end{enumerate}


\end{answer}
